\chapter{Discussion and Conclusions}
\section{Contribution and interpretation}
The project integrates subject-specific delayed entry into the main
single-event survival fitting path of \pkg{}. The default log-rank
statistic and terminal-node estimates consistently use entry-aware risk
sets. Entry information is retained through supported storage, restoration,
resampling, and tuning paths, with explicit errors for unsupported
combinations. The implementation is accompanied by numerical reference
checks and a reproducible simulation framework.

The completed 100-replication study shows lower truth-based survival-curve
error after supplying entry times in all six scenarios for both model
families. Entered-cohort CRPS and AUC error also improve under the chosen
scoring model. Concordance changes are smaller and include reproducible
deterioration for RSF in some settings. Thus, the findings support a
curve-estimation benefit under the studied designs, not a general claim
that delayed-entry adjustment improves every discrimination measure.

\section{Statistical limitations}
The Brier construction is an explicit interpretation of the outline's
observation conditioning, not an equation supplied by the professor.
The added AUC definition also requires scientific review. Their
entered-cohort targets differ from the source-population truth used for
MAD. Pooled censoring weights assume a common censoring hazard; some
simulation mechanisms allow censoring to depend on covariates or entry.
Accordingly, these empirical Brier/AUC comparisons do not prove unbiased
estimation under every simulated censoring mechanism.

Cox uses in-sample curves for outcome-based scores, while RSF uses OOB
curves. That limits cross-family comparisons. All models share preserved
hyperparameters within each fitting comparison, but forests were not
retuned after the package changes. The fixed grids, finite training sample,
and selected DGPs limit generalization. Monte Carlo precision across
100 replications does not remove model misspecification or censoring bias.

The outline also requests fixed-time calibration and survival-curve bias
and variance within selected covariate strata. Those analyses, a real-data
example, and variable-selection performance are not completed by the
present four-metric summary. The literature review and full DGP/source
descriptions also remain to be expanded.

\section{Software scope and computational limitations}
The thesis evaluates single-event log-rank forests with baseline
covariates. Competing-risk LT estimation, time-dependent covariates,
alternative outcome-based split statistics, LT score confidence
intervals, and test-entry-aware native prediction/VIMP remain outside
the established scope. Inclusive entry is deliberate; exact entry/event
ties need explicit interpretation when comparing with strict start--stop
implementations.

All-zero entry times retain the existing fast concordance calculation.
Positive-entry concordance currently compares subject pairs and has
quadratic cost. A fast LT-aware algorithm is a separate future
implementation task; slow pairwise enumeration is not an intrinsic
statistical requirement of left truncation. The local package checks used
the serial configuration, with no sanitizer or OpenMP runtime validation.

\section{Next steps}
The next stage is to review the scoring targets and survival-curve
construction with the professor, complete the requested calibration and
stratum analyses, and expand the literature and DGP descriptions.
Broader external validation and performance work can then support a
future upstream contribution. Upstream code has been integrated locally,
but submission of the LT extension to the original repository has not
been completed. University formatting, authorship details, and final
thesis scope remain to be confirmed.
